% ==========================================================================
% MergeSpec method subsection: "Merge-Aware System Optimization"
% (2026-07-06 REWRITE: the former "Reconstructive Fetching" material was
% removed entirely; this file now holds the system-optimization subsection.
% File name kept so existing \input lines continue to work.)
%
% Insert AFTER \subsection{Relation-Guided Cluster-and-Merge}
% (paper/method_relation_merge.tex).
%
% Notation inherited from the draft and method_relation_merge.tex:
% layer $\ell$, merged $\widehat{\mathbf{W}}_{\ell,j}$ (Eq. 9), groups
% $\mathcal{C}_{\ell,j}$, budget $K$, candidates $\mathcal{M}_\ell$,
% co-occurrence map $\mathbf{C}_{\ell}$, cycle $s$. No new symbols and
% no new equations: this subsection is descriptive system design; a
% schematic figure (TODO: \ref to the hierarchy+pipeline figure) will
% carry the mechanics.
%
% Length target: 0.5--0.75 column-page.
% ==========================================================================

\subsection{Merge-Aware System Optimization}

The draft path of MergeSpec is rebuilt every cycle, so its cost must be
kept off the serving critical path. MergeSpec achieves this with two
system-level mechanisms: a memory hierarchy that decides \emph{where}
merged experts live, and a pipelining schedule that decides \emph{when}
they are built.

\paragraph{Merged-expert memory hierarchy.}
Under uniform coefficients, a merged expert is determined by its member
set alone: $\widehat{\mathbf{W}}_{\ell,j}$ in Eq.~(9) depends on
$\mathcal{C}_{\ell,j}$ but not on the request, the decoding step, or any
routing statistics. Merged experts are therefore content addressable,
and MergeSpec organises them in a two-level hierarchy. The GPU level is
a merged-expert cache: the $K$ resident draft slots at each layer hold
the most recently built merged experts, so the cache lives inside the
draft memory budget of Eq.~(3) and adds nothing to it. The CPU level is
the backing store: it keeps the complete original expert weights, from
which any merged expert can be rebuilt. When the partition at cycle $s$
requests the merged expert of a group, MergeSpec first probes the GPU
cache; a hit returns the slot as is and elides both the member transfers
and the merge. On a miss, the absent members are fetched from CPU
memory and combined on GPU by Eq.~(9), at the cost of one transfer per
absent member.

Cache slots are scarce, so retention follows the probability that a
cached pair is requested again. When the partition changes, MergeSpec
protects first the slots whose two members both remain among the
current candidates $\mathcal{M}_{\ell}$, since the grouping step is
likely to reproduce exactly these pairs, and ranks the remaining slots
by their co-occurrence score; the lowest-ranked slots are evicted.
Both criteria are read directly from state the draft already
maintains. The co-occurrence map thus serves double duty: it guides
which experts merge through Eq.~(8), and it predicts which merged
experts stay useful, so a partition aligned with the temporal locality
of routing converts directly into cache hits and elided merge work.

\paragraph{Fetch-Overlapped Merge Pipelining.}
The second mechanism exploits the fact that CPU--GPU transfers and GPU
computation run on independent engines. In offloaded verification, each
layer must fetch its absent routed experts and execute them. Reading an
expert over the interconnect is one to two orders of magnitude slower
than running it over the speculative window and folding it into a
merge, both of which stream the same weights through GPU memory
bandwidth. Verification is therefore transfer-bound, and the compute
engine is idle for most of each layer. MergeSpec fills this idle time
by scheduling the order in which the layer's experts are served instead
of leaving it arbitrary. Experts that belong to a next-draft group and
are absent from the merged-expert cache are fetched first. Each such
expert enters verification inference as soon as it arrives, and a group
is merged immediately once its last member has been executed, while the
remaining experts of the layer continue streaming behind it. Because
compute is far faster than transfer, the inference and merge work of
the early-scheduled experts completes well before the tail of the
transfer stream, so by the time the layer's last expert arrives, every
merged expert needed by the next draft cycle is already built. Merge
construction and the residual fetches thus hide entirely behind the
transfer bottleneck that offloaded verification must pay anyway, and
the draft phase starts immediately after verification ends. Both
mechanisms leave the draft weights unchanged, so acceptance is
untouched; their effect is purely to remove merge construction from the
critical path.
